\documentclass[11pt,a4paper]{article}
\usepackage[utf8]{inputenc}
\usepackage[french]{babel}
\usepackage[T1]{fontenc}
\usepackage{geometry}
\usepackage{enumitem}
\usepackage{titlesec}
\usepackage{hyperref}
\usepackage{xcolor}
\usepackage{longtable}

\geometry{top=2cm, bottom=2cm, left=2.5cm, right=2.5cm}

\titleformat{\section}
{\normalfont\Large\bfseries\color{blue!70!black}}{\thesection}{1em}{}

\titleformat{\subsection}
{\normalfont\large\bfseries\color{blue!50!black}}{\thesubsection}{1em}{}

\title{\textbf{Guide Développeur 1\\E-Salle ENSAA}}
\author{Académique + Opérations Simples}
\date{\today}

\begin{document}

\maketitle
\tableofcontents
\newpage

\section{Vue d'Ensemble - Développeur 1}

\subsection{Votre Mission}

Vous êtes responsable des \textbf{fondations académiques} et des \textbf{opérations simples} du système E-Salle ENSAA.

\textbf{Vos 4 modules :}
\begin{enumerate}
    \item \textbf{Auth \& Users} (5 jours) - Semaine 1
    \item \textbf{Salle} (6 jours) - Semaine 2
    \item \textbf{Filière} (6 jours) - Semaine 3
    \item \textbf{Réclamation} (6 jours) - Semaine 4
\end{enumerate}

\textbf{Total :} 23 jours de développement pur + 6 jours d'aide à Dev2 (Semaines 5-6)

\subsection{Avantages de Votre Position}

\begin{itemize}
    \item ✅ \textbf{Aucune dépendance externe} : Vous pouvez commencer immédiatement après Semaine 0
    \item ✅ \textbf{Modules complets} : Backend + Frontend pour chaque module
    \item ✅ \textbf{Travail indépendant} : Pas de blocage, pas d'attente
    \item ✅ \textbf{Terminez plus tôt} : Vous aidez Dev2 sur les modules complexes (Emploi, Réservation)
\end{itemize}

\newpage

\section{Vos Modules en Détail}

\subsection{Module 1 : Auth \& Users (Semaine 1 - 5 jours)}

\textbf{Objectif :} Créer le système d'authentification et de gestion des utilisateurs avec table unique et rôles.

\textbf{Entités :}
\begin{itemize}
    \item \texttt{User} : id, nom, prenom, email, password, role (ADMIN/COORDINATEUR/PROFESSEUR/MEMBRE\_CLUB), statut (EN\_ATTENTE/ACTIF/REFUSE), filiere\_id, nom\_club, date\_inscription, date\_approbation, approved\_by
\end{itemize}

\textbf{Fonctionnalités :}
\begin{enumerate}
    \item \textbf{Inscription} : Formulaire (nom, prenom, email, password, role) → Statut = EN\_ATTENTE
    \item \textbf{Login} : Email + Password → Session
    \item \textbf{Compte Admin auto-créé} : Au démarrage de l'application
    \item \textbf{Service Notifications} : Email + WhatsApp (méthode générique)
\end{enumerate}

\textbf{Règles métier :}
\begin{itemize}
    \item Email unique
    \item Password hashé (BCrypt)
    \item Statut par défaut = EN\_ATTENTE
    \item Admin créé automatiquement au premier lancement
\end{itemize}

\textbf{Fichiers à créer :}
\begin{itemize}
    \item \texttt{User.java} (domain)
    \item \texttt{UserRepository.java + UserRepositoryImpl.java}
    \item \texttt{UserService.java + UserServiceImpl.java}
    \item \texttt{AuthServlet.java} (login/logout)
    \item \texttt{RegisterServlet.java} (inscription)
    \item \texttt{NotificationService.java} (Email + WhatsApp)
    \item \texttt{login.jsp}, \texttt{register.jsp}
\end{itemize}

\subsection{Module 2 : Salle (Semaine 2 - 6 jours)}

\textbf{Objectif :} CRUD complet pour la gestion des salles avec types et capacités.

\textbf{Entités :}
\begin{itemize}
    \item \texttt{Salle} : id, nom, type (COURS/TP/TD), capacite, equipements, disponible
\end{itemize}

\textbf{Fonctionnalités :}
\begin{enumerate}
    \item \textbf{CRUD complet} : Créer, Lire, Modifier, Supprimer
    \item \textbf{Filtres} : Par type, par capacité minimale, par disponibilité
    \item \textbf{Recherche} : Par nom
\end{enumerate}

\textbf{Règles métier :}
\begin{itemize}
    \item Nom unique
    \item Capacité > 0
    \item Type obligatoire (COURS, TP, TD)
    \item Salles TP bloquées pour clubs (géré dans Réservation par Dev2)
\end{itemize}

\textbf{Fichiers à créer :}
\begin{itemize}
    \item \texttt{Salle.java} (domain)
    \item \texttt{SalleRepository.java + SalleRepositoryImpl.java}
    \item \texttt{SalleService.java + SalleServiceImpl.java}
    \item \texttt{SalleServlet.java}
    \item \texttt{list.jsp}, \texttt{form.jsp}, \texttt{view.jsp} (dans \texttt{views/salle/})
\end{itemize}

\subsection{Module 3 : Filière (Semaine 3 - 6 jours)}

\textbf{Objectif :} Gestion des filières avec cycles (Préparatoire/Ingénieur) et années.

\textbf{Entités :}
\begin{itemize}
    \item \texttt{Filiere} : id, nom, cycle (PREPARATOIRE/INGENIEUR), annee (1/2/3), effectif, coordinateur\_id
\end{itemize}

\textbf{Fonctionnalités :}
\begin{enumerate}
    \item \textbf{CRUD complet} : Créer, Lire, Modifier, Supprimer
    \item \textbf{Gestion cycles} : Préparatoire (2 ans), Ingénieur (3 ans - DLA1/DLA2/DLA3)
    \item \textbf{Gestion effectifs} : Nombre d'étudiants par filière
    \item \textbf{Assignation coordinateur} : Sélection d'un user avec role COORDINATEUR
\end{enumerate}

\textbf{Règles métier :}
\begin{itemize}
    \item Nom unique
    \item Effectif > 0
    \item Coordinateur doit avoir role = COORDINATEUR
    \item Cycle PREPARATOIRE : années 1-2 uniquement
    \item Cycle INGENIEUR : années 1-3 (DLA1, DLA2, DLA3)
\end{itemize}

\textbf{Fichiers à créer :}
\begin{itemize}
    \item \texttt{Filiere.java} (domain)
    \item \texttt{FiliereRepository.java + FiliereRepositoryImpl.java}
    \item \texttt{FiliereService.java + FiliereServiceImpl.java}
    \item \texttt{FiliereServlet.java}
    \item \texttt{list.jsp}, \texttt{form.jsp}, \texttt{view.jsp} (dans \texttt{views/filiere/})
\end{itemize}

\textbf{⚠️ IMPORTANT :} Dev2 attend que vous terminiez ce module (fin S3) pour commencer Matière !

\subsection{Module 4 : Réclamation (Semaine 4 - 6 jours)}

\textbf{Objectif :} Système de réclamations avec workflow et notifications.

\textbf{Entités :}
\begin{itemize}
    \item \texttt{Reclamation} : id, user\_id, salle\_id, description, urgence (FAIBLE/MOYENNE/ELEVEE), statut (EN\_ATTENTE/TRAITEE), date\_creation, date\_traitement, traite\_par
\end{itemize}

\textbf{Fonctionnalités :}
\begin{enumerate}
    \item \textbf{Créer réclamation} : User sélectionne salle, décrit problème, choisit urgence
    \item \textbf{Notification admin} : Email + WhatsApp à l'admin lors de création
    \item \textbf{Traiter réclamation} : Admin marque comme TRAITEE
    \item \textbf{Notification user} : Email + WhatsApp au user lors du traitement
    \item \textbf{Consulter réclamations} : Liste pour admin et user
\end{enumerate}

\textbf{Règles métier :}
\begin{itemize}
    \item Statut par défaut = EN\_ATTENTE
    \item Notification immédiate à l'admin lors de création
    \item Notification au user lors du traitement
    \item Urgence ELEVEE affichée en rouge dans l'UI
\end{itemize}

\textbf{Fichiers à créer :}
\begin{itemize}
    \item \texttt{Reclamation.java} (domain)
    \item \texttt{ReclamationRepository.java + ReclamationRepositoryImpl.java}
    \item \texttt{ReclamationService.java + ReclamationServiceImpl.java}
    \item \texttt{ReclamationServlet.java}
    \item \texttt{list.jsp}, \texttt{form.jsp}, \texttt{view.jsp} (dans \texttt{views/reclamation/})
\end{itemize}

\newpage

\section{Stratégie Git/GitHub pour Dev1}

\subsection{Configuration Initiale (Semaine 0)}

\textbf{Branches principales :}
\begin{itemize}
    \item \texttt{main} : Production (ne JAMAIS push directement)
    \item \texttt{develop} : Intégration (ne JAMAIS push directement)
    \item \texttt{dev1} : VOTRE branche de travail
\end{itemize}

\textbf{Setup initial :}
\begin{verbatim}
git clone <repo-url>
git checkout -b dev1
git push -u origin dev1
\end{verbatim}

\subsection{Workflow Quotidien}

\textbf{Chaque matin (avant de commencer) :}
\begin{verbatim}
git checkout dev1
git pull origin develop  # Récupérer les changements de Dev2
git pull origin dev1     # Récupérer vos propres changements
\end{verbatim}

\textbf{Pendant le développement :}
\begin{verbatim}
# Commits fréquents (2-3 fois par jour)
git add .
git commit -m "feat(auth): add user registration"
git push origin dev1
\end{verbatim}

\textbf{Conventions de commit :}
\begin{itemize}
    \item \texttt{feat(module): description} - Nouvelle fonctionnalité
    \item \texttt{fix(module): description} - Correction de bug
    \item \texttt{refactor(module): description} - Refactoring
    \item \texttt{docs(module): description} - Documentation
\end{itemize}

\textbf{Exemples pour vos modules :}
\begin{verbatim}
feat(auth): add user login
feat(salle): add CRUD operations
feat(filiere): add cycle management
feat(reclamation): add notification system
fix(auth): fix password hashing
\end{verbatim}

\subsection{Workflow Hebdomadaire}

\textbf{Chaque vendredi (fin de semaine) :}

\textbf{1. Préparer votre code :}
\begin{verbatim}
git checkout dev1
git pull origin develop  # S'assurer d'avoir les derniers changements
# Résoudre conflits si nécessaire
git add .
git commit -m "merge: integrate develop changes"
git push origin dev1
\end{verbatim}

\textbf{2. Créer Pull Request (PR) :}
\begin{itemize}
    \item Aller sur GitHub
    \item Créer PR : \texttt{dev1 → develop}
    \item Titre : \texttt{[S1] Module Auth \& Users - Dev1}
    \item Description : Lister les fonctionnalités ajoutées
    \item Demander review à Dev2
\end{itemize}

\textbf{3. Après approbation :}
\begin{itemize}
    \item Merger la PR sur \texttt{develop}
    \item Dev2 pourra utiliser votre code
\end{itemize}

\subsection{Points de Synchronisation Critiques}

\textbf{Fin Semaine 1 :} Merger Auth \& Users → Dev2 en a besoin pour Administration

\textbf{Fin Semaine 2 :} Merger Salle → Dev2 en a besoin pour Emploi du Temps

\textbf{Fin Semaine 3 :} Merger Filière → Dev2 en a besoin pour Matière ⚠️ \textbf{CRITIQUE !}

\textbf{Fin Semaine 4 :} Merger Réclamation → Vous avez terminé vos modules !

\subsection{Gestion des Conflits}

\textbf{Si conflit lors du pull :}
\begin{verbatim}
git pull origin develop
# Git affiche les conflits
# Ouvrir les fichiers en conflit
# Chercher <<<<<<< HEAD
# Choisir la bonne version ou combiner
git add <fichiers-résolus>
git commit -m "merge: resolve conflicts with develop"
git push origin dev1
\end{verbatim}

\textbf{Conseil :} Communiquez avec Dev2 quotidiennement pour éviter les conflits !

\subsection{Semaines 5-6 : Aide à Dev2}

\textbf{Votre rôle :}
\begin{itemize}
    \item Aider Dev2 sur Emploi du Temps (S5) et Réservation (S6)
    \item Faire des tests sur vos modules
    \item Améliorer l'UI/UX
    \item Corriger les bugs remontés
\end{itemize}

\textbf{Si vous aidez Dev2 sur sa branche :}
\begin{verbatim}
git checkout dev2
git pull origin dev2
# Faire vos modifications
git add .
git commit -m "feat(emploi): help with calendar UI"
git push origin dev2
\end{verbatim}

\newpage

\section{Checklist de Succès}

\subsection{Pour Chaque Module}

\textbf{Avant de merger :}
\begin{itemize}
    \item[$\square$] Entité créée avec annotations Hibernate
    \item[$\square$] Repository implémenté (CRUD complet)
    \item[$\square$] Service implémenté avec logique métier
    \item[$\square$] Servlet créé avec toutes les routes
    \item[$\square$] Pages JSP créées (list, form, view)
    \item[$\square$] Validation des données (Bean Validation)
    \item[$\square$] Gestion des erreurs (try-catch)
    \item[$\square$] Tests manuels effectués
    \item[$\square$] Code compilé sans erreur
    \item[$\square$] Commit + Push sur dev1
    \item[$\square$] PR créée vers develop
\end{itemize}

\subsection{Points de Synchronisation}

\textbf{Fin Semaine 1 :}
\begin{itemize}
    \item[$\square$] Module Auth \& Users terminé et mergé
    \item[$\square$] Dev2 peut commencer Administration
\end{itemize}

\textbf{Fin Semaine 2 :}
\begin{itemize}
    \item[$\square$] Module Salle terminé et mergé
    \item[$\square$] Dev2 peut utiliser Salle dans Emploi du Temps
\end{itemize}

\textbf{Fin Semaine 3 :}
\begin{itemize}
    \item[$\square$] Module Filière terminé et mergé ⚠️ \textbf{CRITIQUE !}
    \item[$\square$] Dev2 peut commencer Matière
\end{itemize}

\textbf{Fin Semaine 4 :}
\begin{itemize}
    \item[$\square$] Module Réclamation terminé et mergé
    \item[$\square$] Tous vos modules sont terminés !
\end{itemize}

\subsection{Communication}

\textbf{Quotidien :}
\begin{itemize}
    \item[$\square$] Standup 10 min avec Dev2 (matin)
    \item[$\square$] Signaler tout blocage immédiatement
\end{itemize}

\textbf{Hebdomadaire :}
\begin{itemize}
    \item[$\square$] Revue de code croisée (vendredi)
    \item[$\square$] PR mergée avant le weekend
\end{itemize}

\section{Conseils Pratiques}

\begin{enumerate}
    \item \textbf{Commencez simple :} CRUD basique d'abord, puis ajoutez la logique métier
    \item \textbf{Testez souvent :} Lancez l'application après chaque fonctionnalité
    \item \textbf{Commitez fréquemment :} 2-3 commits par jour minimum
    \item \textbf{Respectez les deadlines :} Dev2 compte sur vous pour Filière (fin S3) !
    \item \textbf{Communiquez :} Signalez tout problème immédiatement
    \item \textbf{Aidez Dev2 :} Vous terminez plus tôt, profitez-en pour apprendre les modules complexes !
\end{enumerate}

\textbf{Bon développement ! 🚀}

\end{document}

